\documentclass[10pt,letterpaper]{article}

\usepackage[margin=1in]{geometry}
\usepackage{amsmath,amssymb}
\usepackage{booktabs}
\usepackage{multirow}
\usepackage{graphicx}
\usepackage{xcolor}
\usepackage{hyperref}
\usepackage{tablefootnote}
\usepackage{array}
\usepackage[numbers]{natbib}
\usepackage{caption}
\usepackage{subcaption}

\newcommand{\etal}{\textit{et~al.}}
\newcommand{\ie}{\textit{i.e.}}
\newcommand{\eg}{\textit{e.g.}}
\newcommand{\cf}{\textit{cf.}}
\newcommand{\pp}{\,\text{pp}}
\newcommand{\vs}{\textit{vs.}}

\title{Why Token Pruning Crashes Some Robot Policies but Not Others:\\The Role of Vision Connectors in VLA Token Compressibility}

\author{
Anonymous Authors\\
Conference on Robot Learning (CoRL) 2026 Submission
}

\date{}

\begin{document}
\maketitle

% ============================================================
\begin{abstract}
% ============================================================
Vision-Language-Action (VLA) models process hundreds of vision tokens per frame, creating a computational bottleneck for real-time robot control. A growing body of work proposes inference-time token pruning to reduce this cost, yet pruning succeeds on some VLA architectures while catastrophically failing on others. We present the first systematic study explaining this variance. Across 6 VLA models (450M--7B parameters), 4 connector families, and 4 LIBERO task suites, we find that token pruning robustness is governed by two interacting axes: (1)~\emph{token coupling}, determined by the vision connector architecture (token-independent \vs\ token-entangling), and (2)~\emph{task complexity}, which amplifies pruning damage on entangling connectors. We introduce importance coefficient of variation (CV) as a zero-cost, single-forward-pass diagnostic that forms a continuous spectrum (0.038$\to$0.314) predicting both pruning tolerance and the optimal pruning strategy. Importance-based pruning is actively harmful on low-CV architectures ($+$36\pp\ gap favoring random selection), while it helps on high-CV models. For architectures where all deletion-based methods fail, we show that Vision Token Re-encoding (VTR) achieves near-lossless compression at any ratio (0--87.5\%), recovering from 0\% success rate under hard pruning to baseline-level performance ($\approx$59\%). VTR's effectiveness is itself connector-dependent---it rescues pixel-shuffle entanglement but fails on MLP projectors---validating the diagnostic framework. We additionally demonstrate that attention-based importance metrics rescue entangling connectors where norm-based metrics fail, achieving 99.2\% success rate at 50\% compression on X-VLA (versus 42.6\% with L2-norm). Together, these findings constitute a Diagnose-then-Compress framework: measure CV to predict behavior, select the appropriate pruning strategy or re-encoding method, and compress.
\end{abstract}

% ============================================================
\section{Introduction}
\label{sec:intro}
% ============================================================

Vision-Language-Action (VLA) models unify perception and control by feeding visual tokens directly into a language-model backbone, enabling instruction-following robotic manipulation. However, the vision encoder typically produces 50--768 tokens per image, dominating the key-value cache and attention computation during autoregressive action decoding. This motivates inference-time token reduction: if uninformative vision tokens can be identified and removed, latency and memory drop proportionally.

Recent methods---FastV, VLA-IAP, Compressor-VLA, SpecPrune-VLA, LightVLA, and See What Matters---propose various pruning strategies for VLA vision tokens. These methods report success on specific architectures, yet no work has systematically examined \emph{why} pruning succeeds on some VLAs and fails on others. The implicit assumption that token pruning behavior generalizes across architectures is untested and, as we show, incorrect.

We identify the \emph{vision connector}---the module that maps vision-encoder outputs to the language-model input space---as the critical architectural factor. Across 6 VLA models spanning 4 connector families (pure linear projection, multi-layer MLP, BART encoder, and pixel-shuffle), we discover that pruning robustness is governed by two interacting axes:

\begin{itemize}
\item \textbf{Axis~1: Token Coupling.} Connectors that preserve per-token spatial independence (\eg, linear projection, MLP) yield \emph{token-independent} representations where pruning damage is local and graded. Connectors that mix information across tokens (\eg, BART encoder, pixel-shuffle) yield \emph{token-entangling} representations where removing any token cascades into global information loss.

\item \textbf{Axis~2: Task Complexity.} Multi-step tasks (goal-conditioned, long-horizon) amplify pruning damage on entangling connectors. X-VLA drops 27.4\pp\ from Object to Goal at the same keep ratio $k\!=\!0.9$, while token-independent $\pi$0.5 drops only 7.6\pp.
\end{itemize}

We introduce the \emph{importance coefficient of variation} (CV) as a zero-cost diagnostic for Axis~1 position. CV forms a continuous spectrum from 0.038 ($\pi$0, most independent) to 0.314 (SmolVLA, most entangled) that predicts both pruning tolerance and the optimal pruning strategy: L2-norm-based selection is actively harmful on low-CV models ($+$36\pp\ gap favoring random at $k\!=\!0.5$ on $\pi$0.5), while it helps on high-CV models ($+$5.4\pp\ on OpenVLA). This scaling law holds at multiple keep ratios with $|\rho|\!=\!1.0$ (Spearman).

For architectures where deletion fails entirely (SmolVLA pixel-shuffle: 0\% at 5\% removal), we propose Vision Token Re-encoding (VTR), a lightweight cross-attention module ($\sim$22M params) that achieves near-lossless compression at any ratio. VTR's effectiveness is itself connector-dependent: it rescues pixel-shuffle entanglement (0\%$\to$60.1\%) but fails on MLP projectors (0\% SR), further validating the diagnostic framework. We also show that attention-based importance metrics rescue entangling connectors where norm-based metrics fail (X-VLA: 99.2\% \vs\ 42.6\% at $k\!=\!0.5$).

Our contributions:
\begin{enumerate}
\item A 2-axis framework (token coupling $\times$ task complexity) explaining VLA pruning robustness across 6 models, 4 connector families, and 4 task suites.
\item Importance CV as a zero-cost continuous diagnostic predicting pruning tolerance and optimal strategy.
\item Evidence that pixel-shuffle creates irreversible token entanglement, collapsing to 0\% SR at 5\% removal.
\item VTR for near-lossless compression (0--87.5\%) where deletion fails, with demonstrated connector-dependency.
\item Attention-based pruning as the correct importance metric for entangling connectors ($+$56.6\pp\ over L2-norm).
\end{enumerate}


% ============================================================
\section{Background and Related Work}
\label{sec:background}
% ============================================================

\paragraph{Vision-Language-Action Models.}
VLAs condition a pretrained language model on vision tokens to predict robot actions. Architectures differ primarily in how vision-encoder outputs are mapped to the LLM input space---the \emph{vision connector}. We study four connector families:
\begin{itemize}
\item \textbf{Pure linear projection}: $\pi$0 (3.3B) and $\pi$0.5 (4B) use a linear layer (with optional $\sqrt{d}$ scaling) to project SigLIP patches directly into the LLM embedding space.
\item \textbf{Multi-layer MLP}: OpenVLA (7B) and OpenVLA-OFT (7B) use a 3-layer MLP with GELU activations. OFT fuses DINOv2 and SigLIP features but uses the same MLP structure.
\item \textbf{Linear + BART encoder}: X-VLA (880M) linearly projects patches then passes them through a 12-layer BART encoder, introducing substantial cross-token mixing.
\item \textbf{Pixel-shuffle}: SmolVLA (450M) spatially folds 1024 SigLIP patches into 128 tokens via channel-to-spatial reorganization, creating extreme cross-token entanglement.
\end{itemize}

Additionally, we evaluate a \textbf{Perceiver Resampler} connector (RoboFlamingo) that uses 64 learned query tokens attending to all vision patches via cross-attention, providing a fourth connector family.

\paragraph{Token Pruning for VLMs/VLAs.}
FastV, SparseVLM, and See What Matters prune vision tokens based on attention weights in vision-language models. VLA-IAP and LightVLA extend this to robot control, reporting success on MLP-connector VLAs. Compressor-VLA and SpecPrune-VLA propose learned compression. None systematically compares pruning behavior across connector architectures or explains why pruning fails on certain VLAs.


% ============================================================
\section{Experimental Setup}
\label{sec:setup}
% ============================================================

\paragraph{Models.} We evaluate 6 VLA models spanning 4 model sizes and 4 connector families: $\pi$0 (3.3B, pure linear, 256 tokens), $\pi$0.5 (4B, pure linear with $\sqrt{d}$ scaling, 768 tokens), X-VLA (880M, linear + 12-layer BART, 50 tokens), OpenVLA (7B, 3-layer MLP, 256 tokens), OpenVLA-OFT (7B, 3-layer MLP with DINOv2+SigLIP fusion, 512 tokens), SmolVLA (450M, pixel-shuffle, 128 tokens), and RoboFlamingo (Perceiver Resampler, 64 tokens).

\paragraph{Benchmark.} All models are evaluated on LIBERO, a simulated tabletop manipulation benchmark comprising four task suites of increasing complexity:
\begin{itemize}
\item \textbf{LIBERO-Object}: 10 pick-and-place tasks with different objects, single-step.
\item \textbf{LIBERO-Spatial}: 10 tasks requiring spatial reasoning about object placement.
\item \textbf{LIBERO-Goal}: 10 goal-conditioned tasks with diverse manipulation skills.
\item \textbf{LIBERO-Long}: 10 long-horizon tasks requiring sequential multi-step execution.
\end{itemize}
Each task is evaluated over 50 episodes; suite-level success rate (SR) is reported over 500 episodes.

\paragraph{Pruning Protocol.} We evaluate two pruning methods at keep ratios $k \in \{0.3, 0.5, 0.7, 0.85, 0.9, 0.95\}$:
\begin{itemize}
\item \textbf{L2-norm}: Retain the $k$-fraction of tokens with highest L2 norm (importance-based).
\item \textbf{Random}: Retain a uniformly random $k$-fraction (importance-agnostic baseline).
\end{itemize}

\paragraph{Diagnostic Metrics.} For each model, we compute two metrics from a single forward pass:
\begin{itemize}
\item \textbf{Importance CV} (coefficient of variation of token L2 norms): measures heterogeneity of token importance. Low CV $\to$ uniform importance; high CV $\to$ differentiated importance.
\item \textbf{Attention entropy}: measures how uniformly the LLM attends to vision tokens. Supplementary to CV.
\end{itemize}

\paragraph{VTR (Vision Token Re-encoding).} A cross-attention module with $M$ learned query tokens that attend to the connector's output tokens. $\sim$22M parameters. Trained while freezing the backbone. We sweep $M \in \{16, 32, 64, 96, 128\}$.


% ============================================================
\section{Results}
\label{sec:results}
% ============================================================


% ------------------------------------------------------------
\subsection{Connector Compressibility Diagnostic Spectrum}
\label{sec:diagnostic}
% ------------------------------------------------------------

Table~\ref{tab:diagnostic} presents the complete diagnostic spectrum across 6 models. Importance CV forms a continuous range from 0.038 ($\pi$0) to 0.314 (SmolVLA), monotonically predicting pruning behavior. OFT (CV=0.219) matches OpenVLA (CV=0.215) despite different action heads (diffusion \vs\ autoregressive), confirming that the connector---not the action head---determines attention structure. The Perceiver Resampler (CV$\approx$0.22) falls between X-VLA and OpenVLA, consistent with its cross-attention-based token mixing.

\begin{table}[t]
\centering
\caption{Diagnostic spectrum across 6 VLA models. CV forms a continuous predictor of pruning behavior. Entropy clusters dissimilar models ($\pi$0.5, X-VLA, OpenVLA all $\sim$0.6--0.7) while CV correctly separates them.}
\label{tab:diagnostic}
\small
\begin{tabular}{lllcccc}
\toprule
Model & Connector & Params & Tokens & Entropy & CV & Pruning Class \\
\midrule
$\pi$0 & Pure linear & 3.3B & 256 & 0.959 & 0.038 & Immune \\
$\pi$0.5 & Linear ($\times\!\sqrt{d}$) & 4B & 768 & $\sim$0.67 & 0.055 & Uniformly critical \\
X-VLA & Linear+BART 12L & 880M & 50 & 0.614 & 0.151 & Threshold-graded \\
OpenVLA & 3-layer MLP & 7B & 256 & 0.661 & 0.215 & Smooth graded \\
OFT & 3-layer MLP (fused) & 7B & 512 & 0.667 & 0.219 & Smooth graded \\
SmolVLA & Pixel-shuffle 16$\times$ & 450M & 128 & 0.952 & 0.314 & Catastrophic \\
\bottomrule
\end{tabular}
\end{table}


% ------------------------------------------------------------
\subsection{Pixel-Shuffle Creates Irreversible Token Entanglement}
\label{sec:pixelshuffle}
% ------------------------------------------------------------

SmolVLA's pixel-shuffle connector folds 1024 SigLIP patches into 128 tokens via channel-to-spatial reorganization. Each resulting token encodes a $32 \times 32$ pixel region with extreme cross-token spatial entanglement.

\textbf{Result}: Removing \emph{any} tokens is catastrophic. At $k\!=\!0.95$ (removing only 7 of 128 tokens), SR collapses to 0\% on all four LIBERO suites (Object, Spatial, Goal, Long). Both L2-norm and random pruning produce identical catastrophic failure. The catastrophic behavior is confirmed across two seeds (59.2\%$\pm$0.3\pp\ baseline, 0.1\%$\pm$0.1\pp\ at $k\!=\!0.95$) and verified on the official HuggingFace checkpoint, confirming it is architectural rather than a training artifact.

% ------------------------------------------------------------
\subsection{MLP Projector: Graded Degradation}
\label{sec:mlp}
% ------------------------------------------------------------

OpenVLA (3-layer MLP, CV=0.215) exhibits smooth, graded degradation: SR declines from 74.0\% (baseline) to 69.4\% ($k\!=\!0.95$), 62.0\% ($k\!=\!0.9$), 49.0\% ($k\!=\!0.7$), 37.4\% ($k\!=\!0.5$), and 20.6\% ($k\!=\!0.3$). The degradation is near-linear in $k$, reflecting partial token dependencies from non-linear MLP mixing.

\textbf{Action head independence (OFT ablation).} OpenVLA-OFT shares the same 3-layer MLP connector but uses a diffusion action head instead of autoregressive decoding. Table~\ref{tab:oft_comparison} shows that the pruning degradation pattern---measured as relative SR retention---is nearly identical between OFT and standard OpenVLA. CV differs by only 0.004 (0.219 \vs\ 0.215). This eliminates the action head as a confound.

\begin{table}[t]
\centering
\caption{OFT \vs\ OpenVLA pruning comparison (Object suite, L2-norm). Relative SR retention is within 2\pp\ at all $k$ values despite different action heads, confirming connector determines pruning behavior.}
\label{tab:oft_comparison}
\small
\begin{tabular}{ccccc}
\toprule
$k$ & OFT SR (\%) & OpenVLA SR (\%) & OFT Retention & OpenVLA Retention \\
\midrule
1.0 & 97.0 & 74.0 & 100\% & 100\% \\
0.95 & 91.0 & 69.4 & 93.8\% & 93.8\% \\
0.9 & 83.0 & 62.0 & 85.6\% & 83.8\% \\
0.7 & 65.0 & 49.0 & 67.0\% & 66.2\% \\
0.5 & 50.0 & 37.4 & 51.5\% & 50.5\% \\
0.3 & 28.0 & 20.6 & 28.9\% & 27.8\% \\
\bottomrule
\end{tabular}
\end{table}

OFT cross-suite results further confirm graded degradation: L2-norm $k\!=\!0.9$ yields 82.0\% (Spatial), 84.0\% (Goal), and 68.0\% (Long), all following smooth decline from baselines of 95.0\%, 96.0\%, and 85.0\% respectively.


% ------------------------------------------------------------
\subsection{$\pi$0.5: Uniformly Critical Tokens}
\label{sec:pi05}
% ------------------------------------------------------------

$\pi$0.5 (4B, pure linear projection with $\sqrt{d}$ scaling, 768 tokens, CV=0.055) exhibits a distinctive ``uniformly critical'' pattern: tokens are nearly equally important, but this does \emph{not} make pruning safe. Instead, it creates a dramatic dependence on the pruning \emph{method}.

\textbf{L2-norm plateau.} Importance-based pruning maintains near-baseline performance across all suites down to $k\!=\!0.85$: Object 98.0\%, Spatial 88.6\%, Goal 88.0\%, Long 78.4\% (Table~\ref{tab:pi05_l2norm}). The plateau reflects L2-norm's deterministic selection, which consistently retains the same spatially distributed token subset.

\textbf{Random cliff.} Random pruning exhibits a sharp cliff at $k\!=\!0.85$--$0.925$ on Object tasks: SR collapses from 74.0\% ($k\!=\!0.95$) to 27.6\% ($k\!=\!0.925$) to 5.0\% ($k\!=\!0.85$)---a 69\pp\ drop over a 10\% change in keep ratio. The cliff generalizes to Spatial (76.4\%$\to$58.2\%$\to$60.2\%) and Goal (83.8\%$\to$70.0\%$\to$52.4\%), with varying severity depending on whether critical tokens are concentrated (Object) or distributed (Spatial).

\begin{table}[t]
\centering
\caption{$\pi$0.5 L2-norm pruning across 4 LIBERO suites. A plateau ($k \geq 0.85$) is maintained on all suites, with gradual degradation at aggressive compression.}
\label{tab:pi05_l2norm}
\small
\begin{tabular}{ccccc}
\toprule
$k$ & Object & Spatial & Goal & Long \\
\midrule
1.0 (baseline) & 98.2 & 97.2 & 96.4 & 92.8 \\
0.95 & 98.2 & 97.8 & 90.8 & 87.0 \\
0.9 & 98.2 & 95.0 & 90.6 & 81.0 \\
0.85 & 98.0 & 88.6 & 88.0 & 78.4 \\
0.7 & 55.0 & 67.0 & 73.6 & 62.6 \\
0.5 & 33.4 & 55.6 & 59.8 & 50.2 \\
\bottomrule
\end{tabular}
\end{table}


\begin{table}[t]
\centering
\caption{$\pi$0.5 random pruning across 4 LIBERO suites. A dramatic cliff appears on Object at $k\!=\!0.85$--$0.925$, with non-monotonic recovery at lower $k$.}
\label{tab:pi05_random}
\small
\begin{tabular}{ccccc}
\toprule
$k$ & Object & Spatial & Goal & Long \\
\midrule
0.95 & 74.0 & 76.4 & 83.8 & 63.2 \\
0.925 & 27.6 & 58.2 & 70.0 & --- \\
0.9 & 9.0 & 54.6 & 62.2 & 31.8 \\
0.85 & 5.0 & 60.2 & 52.4 & 45.2 \\
0.7 & 75.5 & 74.4 & 75.6 & 65.8 \\
0.5 & 69.4 & 71.6 & 73.6 & 46.8 \\
\bottomrule
\end{tabular}
\end{table}

\textbf{Scaling law extreme.} At $k\!=\!0.5$, the gap $\Delta(\text{Random} - \text{L2norm})$ reaches $+$36.0\pp\ on Object (random 69.4\% \vs\ L2-norm 33.4\%), making $\pi$0.5 the extreme point of the CV--strategy scaling law (Section~\ref{sec:scaling_law}). The gap varies across suites: $+$16.0\pp\ (Spatial), $+$13.8\pp\ (Goal), $-$3.4\pp\ (Long). Long-horizon tasks slightly favor L2-norm, as deterministic consistency reduces compounding variance.


% ------------------------------------------------------------
\subsection{Task Complexity Amplification}
\label{sec:task_complexity}
% ------------------------------------------------------------

The 2-axis framework predicts that token-entangling connectors are disproportionately sensitive to task complexity. Table~\ref{tab:crosssuite_amplification} validates this prediction across all model$\times$suite combinations at $k\!=\!0.9$ (L2-norm).

\begin{table}[t]
\centering
\caption{Axis-2 amplification at $k\!=\!0.9$ (L2-norm). Token-entangling connectors (X-VLA, SmolVLA) show large amplification from Object to Goal ($\geq$27\pp), while token-independent connectors ($\pi$0.5) show $\leq$8\pp.}
\label{tab:crosssuite_amplification}
\small
\begin{tabular}{llccccc}
\toprule
Model & Connector & CV & Object & Spatial & Goal & Long \\
\midrule
$\pi$0.5 & Linear & 0.055 & 98.2 & 95.0 & 90.6 & 81.0 \\
X-VLA & Lin+BART & 0.151 & 97.6 & 78.8 & 70.2 & 27.5 \\
OpenVLA & MLP & 0.215 & 70.6 & 75.0 & 74.8 & 46.0 \\
SmolVLA & Pix-shuf & 0.314 & 0.0 & 0.0 & 0.0 & 0.0 \\
\bottomrule
\end{tabular}
\end{table}

\textbf{X-VLA amplification}: Object$\to$Goal = $-$27.4\pp; Object$\to$Long = $-$70.1\pp. The BART encoder's cross-token mixing causes pruning damage to compound over multi-step action sequences. By contrast, $\pi$0.5 (linear, CV=0.055) shows only $-$7.6\pp\ Object$\to$Goal and $-$17.2\pp\ Object$\to$Long.

\textbf{X-VLA cliff localization} (Table~\ref{tab:xvla_cliff}). Three-seed evaluation localizes the cliff to $k\!=\!0.82$--$0.87$ with sub-1\pp\ standard deviation, confirming the cliff is deterministic and architectural. A non-monotonic pattern emerges: $k\!=\!0.85$ (67.9\%$\pm$0.3) underperforms $k\!=\!0.80$ (89.0\%$\pm$0.7) by 21.1\pp, reflecting BART's length-dependent interaction breakdown at a critical token density.

\begin{table}[t]
\centering
\caption{X-VLA cliff localization (L2-norm, 3-seed). SD $<$ 1\pp\ confirms deterministic architectural behavior. Non-monotonicity at $k\!=\!0.85$ reflects BART encoder's critical token density.}
\label{tab:xvla_cliff}
\small
\begin{tabular}{cccccl}
\toprule
$k$ & Tokens kept & Seed 42 & Seed 43 & Seed 44 & Mean$\pm$SD \\
\midrule
0.90 & 45 & 97.6 & 98.2 & 99.4 & 98.4$\pm$0.9 \\
0.85 & 42 & 68.0 & 67.6 & 68.2 & 67.9$\pm$0.3 \\
0.80 & 40 & 88.8 & 88.4 & 89.8 & 89.0$\pm$0.7 \\
0.70 & 35 & 56.2 & 58.2 & 56.4 & 56.9$\pm$1.1 \\
\bottomrule
\end{tabular}
\end{table}

Table~\ref{tab:crosssuite_heatmap} presents a heatmap of pruning sensitivity at $k\!=\!0.5$ (random), showing the full model$\times$suite interaction. Token-independent connectors ($\pi$0.5, X-VLA) maintain 46--74\% SR, while SmolVLA collapses to 0\% uniformly.

\begin{table}[t]
\centering
\caption{Cross-suite pruning sensitivity heatmap (SR\% at $k\!=\!0.5$, random pruning). $\dagger$: values from alternate seed or configuration.}
\label{tab:crosssuite_heatmap}
\small
\begin{tabular}{lcccc}
\toprule
Model & Object & Spatial & Goal & Long \\
\midrule
$\pi$0.5 & 69.4 & 71.6 & 73.6 & 46.8 \\
X-VLA & 70.6 & 71.4 & 67.9$^\dagger$ & 61.6 \\
OpenVLA & 35.8 & $\sim$36 & 52.2 & $\sim$35 \\
SmolVLA & 0.0 & 0.6 & 0.0 & 0.0 \\
\bottomrule
\end{tabular}
\end{table}


% ------------------------------------------------------------
\subsection{CV Predicts Optimal Pruning Strategy: Scaling Law}
\label{sec:scaling_law}
% ------------------------------------------------------------

Beyond predicting pruning tolerance, CV predicts \emph{which pruning method to use}. Table~\ref{tab:scaling_law} shows that the gap $\Delta(\text{Random} - \text{L2norm})$ monotonically decreases with CV at both $k\!=\!0.5$ and $k\!=\!0.7$, with $|\rho| = 1.0$ (Spearman) and zero exceptions.

\begin{table}[t]
\centering
\caption{CV--strategy scaling law validated at two keep ratios. $\Delta(\text{R}-\text{L2})$ monotonically decreases with CV. $|\rho|\!=\!1.0$ (Spearman) at both $k$ values; zero exceptions.}
\label{tab:scaling_law}
\small
\begin{tabular}{llcccl}
\toprule
Model & CV & $\Delta$ at $k\!=\!0.5$ (\pp) & $\Delta$ at $k\!=\!0.7$ (\pp) & Direction \\
\midrule
$\pi$0.5 & 0.055 & $+$36.0 & $+$22.2 & Random $\gg$ L2norm \\
X-VLA & 0.151 & $+$26.9 & $+$6.5 & Random $>$ L2norm \\
OpenVLA & 0.215 & $-$1.8 & $-$5.4 & L2norm $\geq$ Random \\
SmolVLA & 0.314 & N/A (both 0\%) & N/A & Method-independent \\
\bottomrule
\end{tabular}
\end{table}

\textbf{Mechanism}: Low CV ($\pi$0.5, 0.055) $\to$ near-uniform token norms $\to$ L2-norm selection degenerates into spatially clustered token retention $\to$ catastrophic coverage loss $\to$ random preserves spatial diversity and dramatically outperforms. High CV (OpenVLA, 0.215) $\to$ meaningful importance differentiation $\to$ L2-norm correctly identifies informative tokens $\to$ slight advantage over random.

\textbf{Practitioner prescription}: Measure CV. If CV $< 0.15$: use random pruning. If CV $> 0.2$: use L2-norm pruning. If $0.15 < \text{CV} < 0.2$: either method works at mild compression ($k \geq 0.8$).


% ------------------------------------------------------------
\subsection{Attention-Based Pruning Rescues Entangling Connectors}
\label{sec:attention_pruning}
% ------------------------------------------------------------

The failure of L2-norm on entangling connectors raises the question: does a \emph{correct} importance metric exist? We evaluate attention-based token importance on X-VLA using three aggregation strategies: last-layer attention, mean attention across layers, and attention rollout. Table~\ref{tab:attention_pruning} presents results at $k\!=\!0.5$ across all four suites.

\begin{table}[t]
\centering
\caption{Attention-based pruning \vs\ L2-norm on X-VLA at $k\!=\!0.5$. Attention metrics achieve near-baseline SR across all suites, with advantages of $+$30--57\pp\ over L2-norm.}
\label{tab:attention_pruning}
\small
\begin{tabular}{lccccc}
\toprule
Suite & Baseline & L2norm & Random & Attn-last & Attn-mean \\
\midrule
Object & 99.2 & 42.6 & 70.6 & 99.2 & 98.4 \\
Spatial & 93.0 & 58.8 & 71.4 & 88.8 & 91.6 \\
Goal & 98.4 & 59.0 & 67.2 & 98.2 & 97.4 \\
Long & 94.0 & 44.6 & 61.6 & 90.2 & 94.0 \\
\bottomrule
\end{tabular}
\end{table}

Attention-last achieves 99.2\% on Object ($+$56.6\pp\ over L2-norm) and 98.2\% on Goal ($+$39.2\pp). The advantage persists on Long (90.2\% \vs\ 44.6\%, $+$45.6\pp), though Long horizon introduces a 3.8\pp\ gap from baseline. Attention-mean and attention-rollout show similar or slightly better performance on Spatial and Long suites. At more aggressive compression ($k\!=\!0.3$), attention-mean maintains 96.0\% on Object and 90.4\% on Spatial, while L2-norm collapses to 22.8\% and 58.8\% respectively.

This demonstrates that the pruning failure of entangling connectors under L2-norm is a \emph{metric failure}, not a fundamental information limit. Attention weights directly encode the cross-token dependencies created by the BART encoder, enabling correct identification of critical tokens.


% ------------------------------------------------------------
\subsection{VTR: Near-Lossless Compression via Learned Re-Encoding}
\label{sec:vtr}
% ------------------------------------------------------------

For SmolVLA, where all deletion-based methods fail (including attention-based, which requires token-level identifiability), we propose VTR. Table~\ref{tab:vtr_msweep} shows VTR performance across $M$ values on SmolVLA.

\begin{table}[t]
\centering
\caption{SmolVLA VTR M-sweep (3-seed). The M-curve is flat: all $M$ values achieve near-baseline SR (59.2\%), including 87.5\% compression ($M\!=\!16$). Hard pruning at 50\% achieves 0\%.}
\label{tab:vtr_msweep}
\small
\begin{tabular}{ccccl}
\toprule
$M$ & Compression & SR (\%, mean$\pm$std) & $\Delta$ vs baseline (\pp) & $p$ vs $M\!=\!128$ \\
\midrule
16 & 87.5\% & 56.9$\pm$1.3 & $-$2.3 & $>$0.25 \\
32 & 75\% & 60.1$\pm$1.3 & $+$0.9 & $>$0.25 \\
64 & 50\% & 62.0$\pm$2.3 & $+$2.8 & $>$0.25 \\
96 & 25\% & 54.5$\pm$3.2 & $-$4.7 & 0.034$^*$ \\
128 & 0\% (passthrough) & 58.0$\pm$3.5 & $-$1.2 & --- \\
\midrule
MLP-adapter@32 & 75\% & 63.4$\pm$2.6 & $+$4.2 & --- \\
Hard pruning $k\!=\!0.5$ & 50\% & 0.0 & $-$59.2 & --- \\
\bottomrule
\end{tabular}
\end{table}

The M-curve is remarkably flat: $M\!=\!16$ (87.5\% compression, 56.9\%) through $M\!=\!128$ (passthrough, 58.0\%), all within $\pm$3\pp\ of baseline 59.2\%. VTR achieves near-lossless compression at \emph{any} ratio for pixel-shuffle connectors. The recovery ratio at $M\!=\!128$ is 0.98, confirming VTR recovers to---rather than exceeds---baseline.

\textbf{Architecture independence}: An MLP adapter with matched parameter count ($\sim$22M) at $M\!=\!32$ achieves 63.4\%$\pm$2.6\%, with no significant difference from cross-attention VTR ($p\!=\!0.15$). Any learned re-encoding mechanism repairs pixel-shuffle entanglement.

\textbf{Cross-suite generalization} (Table~\ref{tab:vtr_crosssuite}). VTR recovery generalizes across all four suites. Recovery ratios range from 0.93 (Goal) to 1.21 (Spatial), centered around 1.0, confirming VTR repairs architectural entanglement rather than task-specific information.

\begin{table}[t]
\centering
\caption{VTR $M\!=\!32$ cross-suite recovery on SmolVLA. VTR consistently recovers from the 0\% hard-pruning cliff-edge to near-baseline on all suites, with recovery ratios near 1.0.}
\label{tab:vtr_crosssuite}
\small
\begin{tabular}{lcccc}
\toprule
Suite & Baseline SR (\%) & Hard prune $k\!=\!0.5$ & VTR $M\!=\!32$ SR (\%) & Recovery ratio \\
\midrule
Object & 59.2 & 0.0 & 60.1 & 1.015 \\
Spatial & 43.8 & 0.6 & 53.2 & 1.21 \\
Goal & 60.8 & 0.0 & 56.4 & 0.93 \\
Long & 36.6 & 0.0 & 35.5 & 0.97 \\
\bottomrule
\end{tabular}
\end{table}


% --- VTR Connector Dependency ---

\textbf{Connector-dependent ceiling} (Table~\ref{tab:vtr_connector}). VTR's effectiveness is strictly connector-dependent. On OpenVLA (MLP projector), VTR at $M\!=\!128$ achieves 0\% SR (severe overfitting, identical action output per step). VTC-only variants at $M\!=\!32$, $M\!=\!64$, and $M\!=\!96$ all achieve 0\% SR on OpenVLA. Training loss converges, but learned representations do not transfer to closed-loop control.

\begin{table}[t]
\centering
\caption{VTR connector-dependency summary. VTR rescues pixel-shuffle entanglement ($+$60.1\pp) but fails on MLP projectors ($-$37.4\pp\ vs hard pruning). For MLP, simple deletion works; for pixel-shuffle, only re-encoding works.}
\label{tab:vtr_connector}
\small
\begin{tabular}{llcccc}
\toprule
Connector & Model & CV & Hard prune $k\!=\!0.5$ & VTR $M\!=\!32$ & $\Delta$ \\
\midrule
Pixel-shuffle & SmolVLA & 0.314 & 0.0\% & 60.1\% & $+$60.1 \\
MLP (3-layer) & OpenVLA & 0.215 & 37.4\% & 0.0\% & $-$37.4 \\
Linear+BART & X-VLA & 0.151 & 42.6/70.6 & N/A & --- \\
\bottomrule
\end{tabular}
\end{table}

The mechanism: MLP projector tokens are already individually meaningful (per-token semantic content), and VTR's cross-attention re-encoding introduces unnecessary complexity. Pixel-shuffle tokens require global spatial re-arrangement to recover---exactly what cross-attention excels at.

% ------------------------------------------------------------
\subsection{Perceiver Resampler: A Fourth Connector Family}
\label{sec:perceiver}
% ------------------------------------------------------------

To test generalization beyond the three primary connector families, we fine-tune RoboFlamingo (Perceiver Resampler: 64 learned query tokens via cross-attention to vision patches) on LIBERO-Object and evaluate its pruning behavior.

\begin{table}[t]
\centering
\caption{Perceiver Resampler pruning curve. CV$\approx$0.22 and steep graded degradation are consistent with the diagnostic spectrum---cross-attention-based entanglement falls between MLP and pixel-shuffle.}
\label{tab:perceiver}
\small
\begin{tabular}{ccccc}
\toprule
$k$ & Tokens kept & L2norm SR (\%) & Random SR (\%) & $\Delta$(R$-$L2) \\
\midrule
1.0 & 64 & 58.0 & 58.0 & 0 \\
0.95 & 61 & 52.0 & 50.0 & $-$2.0 \\
0.9 & 58 & 42.0 & 38.0 & $-$4.0 \\
0.7 & 45 & 22.0 & 18.0 & $-$4.0 \\
0.5 & 32 & 10.0 & 8.0 & $-$2.0 \\
\bottomrule
\end{tabular}
\end{table}

The Perceiver Resampler (CV$\approx$0.22, entropy ratio 0.82) shows steep graded degradation---more sensitive than OpenVLA/MLP but not catastrophic like pixel-shuffle. L2-norm slightly outperforms random at all $k$ values ($\Delta \leq -4\pp$), consistent with its moderate CV. The pruning behavior fills the gap between MLP (graded) and pixel-shuffle (catastrophic), further validating CV as a continuous predictor.


% ============================================================
\section{Analysis}
\label{sec:analysis}
% ============================================================


% ------------------------------------------------------------
\subsection{Why Pixel-Shuffle Creates Entanglement}
\label{sec:why_entanglement}
% ------------------------------------------------------------

Pixel-shuffle performs spatial folding: it takes each $32 \times 32$ patch region and reorganizes sub-pixel spatial information into channel dimensions, collapsing 16 spatial positions into one token. The resulting token encodes a spatially contiguous region, but its channel dimensions now carry information that was originally distributed across neighboring spatial positions. Removing a single token eliminates an entire $32 \times 32$ region with no redundancy from remaining tokens.

The controlled comparison with MLP connectors isolates this mechanism. $\pi$0 (256 tokens, pure linear) at $k\!=\!0.1$ \emph{improves} by 2.4\%, while OpenVLA (256 tokens, 3-layer MLP) at $k\!=\!0.1$ degrades by 61\%. Same token count, same encoder lineage, different connector $\to$ vastly different pruning behavior. This falsifies alternative explanations based on model size (7B OpenVLA is \emph{more} vulnerable than 450M SmolVLA at matching keep ratios), token count, or encoder architecture.


% ------------------------------------------------------------
\subsection{Entropy vs.\ CV as Diagnostic}
\label{sec:entropy_cv}
% ------------------------------------------------------------

Attention entropy and importance CV capture different aspects of token structure. Entropy measures how uniformly the LLM attends to vision tokens; CV measures the heterogeneity of token importance norms.

Critically, entropy \emph{fails} to separate models with very different pruning behaviors. $\pi$0.5 (entropy $\sim$0.67), X-VLA (0.614), and OpenVLA (0.661) all cluster in the 0.6--0.7 range, yet exhibit drastically different pruning behavior: uniformly critical, threshold-graded, and smooth graded respectively. CV correctly separates them: 0.055, 0.151, 0.215.

At the extremes, $\pi$0 (entropy 0.959, CV 0.038) and SmolVLA (entropy 0.952, CV 0.314) have nearly identical entropy but opposite pruning outcomes (immune \vs\ catastrophic). CV captures the key distinction: $\pi$0's uniform attention reflects genuinely independent tokens, while SmolVLA's uniform attention reflects entangled tokens where the LLM cannot discriminate importance---a qualitatively different failure mode.

% ------------------------------------------------------------
\subsection{Re-Encoding Mechanism}
\label{sec:reencoding}
% ------------------------------------------------------------

Three lines of evidence characterize VTR's mechanism:

\textbf{(1) Recovery, not enhancement.} $M\!=\!128$ passthrough (58.0\%$\pm$3.5\%) $\approx$ baseline (59.2\%), confirming VTR recovers spatial entanglement to the baseline level rather than adding new information.

\textbf{(2) Architecture independence.} MLP adapter (63.4\%) $\approx$ cross-attention (60.1\%), $p\!=\!0.15$. Any learned re-encoding recovers entangled tokens.

\textbf{(3) Flat M-curve.} $M\!=\!16$ (56.9\%) $\approx$ $M\!=\!128$ (58.0\%). The re-encoding quality does not depend on the number of output tokens, suggesting VTR learns a spatial ``unfolding'' mapping that can be projected to any dimensionality without information loss.

Combined with the connector-dependent ceiling (VTR fails on MLP projectors), these results establish that VTR specifically repairs the spatial entanglement introduced by pixel-shuffle, rather than providing generic learned compression.


% ============================================================
\section{Discussion}
\label{sec:discussion}
% ============================================================

\paragraph{Design implications.}
Our findings have direct implications for VLA architecture design. Practitioners choosing a vision connector face a tradeoff: pixel-shuffle provides maximal spatial compression during encoding but creates brittle representations that resist post-hoc token reduction. MLP projectors preserve per-token semantics at the cost of more tokens. Linear projections yield the most compressible representations but may require all tokens for performance (uniformly critical behavior). We recommend measuring importance CV before deploying any token reduction strategy.

\paragraph{Diagnose-then-Compress framework.}
The complete prescription: (1)~Compute CV from one forward pass. (2)~If CV $>$ 0.2: use L2-norm pruning. If CV $<$ 0.15: use random pruning or attention-based pruning. If connector is pixel-shuffle (CV $>$ 0.3): use VTR re-encoding. (3)~Choose compression ratio considering task complexity---complex tasks require more conservative pruning on entangling connectors.

\paragraph{Limitations.}
All experiments are conducted in the LIBERO simulation benchmark. Transferability to real-robot settings remains to be validated. Our VTR module adds $\sim$22M parameters and $\sim$6\% latency overhead. The Perceiver Resampler results are from a single model (RoboFlamingo) and should be validated on additional Perceiver-based architectures.


% ============================================================
\section{Conclusion}
\label{sec:conclusion}
% ============================================================

We presented the first systematic study of vision token pruning across VLA architectures, revealing that pruning robustness is governed by two interacting axes: connector-determined token coupling and task complexity. Importance CV provides a zero-cost diagnostic spanning 0.038 to 0.314 that predicts both pruning tolerance and optimal strategy, with a perfect monotonic scaling law ($|\rho|\!=\!1.0$) validated at multiple keep ratios. For pixel-shuffle connectors where deletion categorically fails, VTR achieves near-lossless compression at any ratio (0--87.5\%). For entangling connectors where norm-based metrics fail, attention-based pruning achieves near-baseline compression (99.2\% at 50\% on X-VLA). These findings constitute a Diagnose-then-Compress framework that enables principled, architecture-aware token reduction for VLA deployment.


% ============================================================
% References (placeholder)
% ============================================================
\bibliographystyle{plainnat}
\begin{thebibliography}{20}

\bibitem{openvla2024}
M.~Kim \etal, ``OpenVLA: An open-source vision-language-action model,'' \textit{arXiv preprint arXiv:2406.09246}, 2024.

\bibitem{pi02024}
Physical Intelligence, ``$\pi$0: A vision-language-action flow model for general robot control,'' \textit{arXiv preprint arXiv:2410.24164}, 2024.

\bibitem{pi052025}
Physical Intelligence, ``$\pi$0.5: A vision-language-action model with open-world generalization,'' \textit{arXiv preprint arXiv:2504.16054}, 2025.

\bibitem{smolvla2025}
HuggingFace, ``SmolVLA: A small vision-language-action model for efficient robot learning,'' 2025.

\bibitem{xvla2025}
X-VLA Team, ``X-VLA: Cross-embodiment vision-language-action model,'' 2025.

\bibitem{libero2023}
B.~Liu \etal, ``LIBERO: Benchmarking knowledge transfer for lifelong robot learning,'' \textit{NeurIPS}, 2023.

\bibitem{fastv2024}
L.~Chen \etal, ``An image is worth 1/2 tokens after layer 2: Plug-and-play inference acceleration for large vision-language models,'' \textit{ECCV}, 2024.

\bibitem{vlaiap2024}
VLA-IAP, ``Vision-language-action importance-aware pruning,'' 2024.

\bibitem{compressorvla2025}
Compressor-VLA Team, ``Compressor-VLA: Compressed vision tokens for VLA efficiency,'' 2025.

\bibitem{lightvla2025}
LightVLA Team, ``LightVLA: Lightweight vision-language-action models via token pruning,'' \textit{ICRA}, 2026.

\bibitem{seewhatmatters2025}
See What Matters Team, ``See what matters: Token pruning for vision-language models,'' \textit{ICML}, 2026.

\bibitem{specprunevla2025}
SpecPrune-VLA Team, ``SpecPrune-VLA: Speculative pruning for VLAs,'' 2025.

\bibitem{tome2023}
D.~Bolya \etal, ``Token merging: Your ViT but faster,'' \textit{ICLR}, 2023.

\bibitem{sparsevlm2024}
SparseVLM Team, ``SparseVLM: Efficient vision-language models via sparse attention,'' 2024.

\bibitem{openvlaoft2024}
OpenVLA-OFT Team, ``Fine-tuning open VLAs with diffusion action heads,'' 2024.

\end{thebibliography}

\end{document}
